\chapter{Suchbäume und Grundlagen}\label{ch:foundations}

% [SKELETT] REIN KONZEPTIONELL + VISUELL, KEINE Katalog-Tabellen. Dreischritt Konzept->Hardware->Software:
% Abstraktions-Verantwortlichkeiten in 2.1 konzeptionell verorten, in 2.3 mit den statischen Hardware-/
% Compile-Time-Mitteln aus 2.2 umsetzen. Alle Tabellen liegen NICHT hier, sondern in Kapitel 3.

Dieses Kapitel führt vom vertrauten Gegenstand --- den Suchbäumen --- über die Hardware, die ihre Leistung
prägt, zu den Software-Mitteln, mit denen sich ihre Entwurfsentscheidungen abstrahieren lassen. Alles bleibt
\emph{bildhaft und konzeptionell}; die systematischen Kataloge folgen in Kapitel~\ref{ch:measurement}.

\section{Konzept und Überblick über den Stand der Technik}\label{sec:overview}
Wir kartieren zunächst die Landschaft der Suchstrukturen --- vergleichsbasierte Bäume, digitale und
präfixbasierte Verfahren, Hashing, räumliche und flache Strukturen --- und stellen \emph{jeden} bekannten
Algorithmus bildhaft im Detail dar. Hier wird die Abstraktion bereits konzeptionell verortet: jeder Algorithmus
ist ein Punkt in einem gemeinsamen Entwurfsraum.
\bildPH{B3: Such-Struktur-Landkarte}{Taxonomie: vergleichsbasiert / digital / Hash / räumlich / flach.}
\bildPH{B5: Entwurfsraum (Konzept)}{Achsen-Würfel; ein Stand-der-Technik-Algorithmus = ein Punkt.}
\inhaltPH{Landschaft aus 02\_fundamentals (sec:search-classes) + SOTA-Überblick aus 03\_state\_of\_the\_art (sec:sota-overview). \textbf{Bildserie SOTA-ALGO-1..M:} je Paper-Algorithmus (ART, HOT, START, SuRF, Wormhole, B$^2$-Baum, \ldots) ein Detailbild.}

\section{Hardware: historische Anpassung an die Suchleistung}\label{sec:hardware-history}
Die Leistung von Suchstrukturen wird seit jeher von der Hardware getrieben. Wir führen die historische
Einführung cache-naher Hardware ein --- Cache-Hierarchie, Cache-Line, Adressübersetzung, Kohärenz --- mit Fokus
auf die \emph{statischen und übersetzungszeit-relevanten} Eigenschaften, die später die Abstraktion tragen.
\bildPH{B4: Cache-Hierarchie und Cache-Line}{Speicher-Pyramide + eine Cache-Line mit belegten/leeren Bytes (Auslastung).}
\inhaltPH{Hardware aus 02\_fundamentals (sec:cache-basics) + 03\_state\_of\_the\_art (sec:sota-cache)}

\section{Software-Mittel und Zusammenführung zur Abstraktion}\label{sec:software-means}
Mit welchen Software-System-Mitteln --- Übersetzungszeit-Metaprogrammierung, statische Nutzung der Hardware aus
Abschnitt~\ref{sec:hardware-history} --- lässt sich die Abstraktion \emph{umsetzen}? Hier führen wir die bekannten
Verfahren bildhaft zu \emph{unseren} Achsen zusammen und definieren alle Bestandteile, die konzeptionellen
Spezialkonstrukte (Hüllen, besondere Entwurfsmuster), die Begrifflichkeiten und die Grundlagen wissenschaftlichen
Messens.
\bildPH{ANAT-1..3: Anatomie-Brücke}{Menschliche Anatomie (Organe + Verdrahtung) $\to$ Achse=Organ $\to$ Mess-System als Diagnostik.}
\bildPH{SYNTH: Synthese Stand der Technik $\to$ Achsen}{Bausteine der Paper-Algorithmen fließen in die gemeinsame Achsen-Matrix.}
\bildPH{B6: Achse = Organ, Anatomie = Verdrahtung}{Organe + generalisierte Interfaces dazwischen (z.\,B. das Allokations-Interface).}
\bildPH{B7: Drei Ebenen und Gattungen}{Wurzel $\to$ Gattungen $\to$ Tier-Unterklassen $\to$ Organe.}
\bildPH{B8: Hüllen / ABI-Adapter}{Lebewesen in der ABI-Hülle über der Modulgrenze; die SearchEngine-Sicht.}
\inhaltPH{Definitionen aus 04\_concept\_architecture (Eine-Architektur/3-Ebenen/ABI) + 02\_fundamentals (Vergleichsinterfaces, C++23-Pattern, wissenschaftliches Messen) + 03 (Design-Pattern-Beitrag). \textbf{Weitere Bilder:} GATT (alle Gattungen), PATTERN (Entwurfsmuster), USAGE (welche Achse nutzt welche), UML-* (Landkarten aller Code-Interfaces).}
